% --------------------- NO TOCAR ----------------------
\documentclass[a4paper, 10pt, twocolumn]{report}
\input plantillaTeXpracticasRL.sty % Incluir fichero de estilo del máster
% -----------------------------------------------------


\renewcommand{\epigraphflush}{flushleft}

% MODIFICA ESTOS DATOS
% Incluir el nombre de la obra y los nombres completos de los autores del trabajo, junto con el correo electrónico um.es para cada autor.
\title{\textbf{\textcolor{UMUblue}{Extensiones de Machine Learning \\ Aprendizaje por Refuerzo}}}
\author{
Marta Lucas Marín 	\\
\texttt{m.lucasmarin@um.es} \\[0.4cm]
Francisco José Lucas Robles 	\\
\texttt{fjlr1@um.es} \\[0.4cm]
Jorge Ruiz Ortiz 	\\
\texttt{j.ruizortiz@um.es}}
\date{
\today
\vspace{3cm}%  
\epigraph{\large{¡Nunca seré derrotado por una máquina! Jamás se inventará un programa que supere la inteligencia humana.}}{\large{G. Kasparov, 1989}\footnote{De una entrevista publicada en la edición nº 55 de la revista \textit{Jeux \& Stratégie.}}}
}

\setcounter{tocdepth}{1} % Profundidad del índice

\begin{document}

% Carátula
\begin{titlepage}
\newgeometry{margin=1in}
\maketitle
\restoregeometry
\end{titlepage}

% Resumen
% ============================================================
%  Resumen — Extensiones de Machine Learning: Aprendizaje por Refuerzo
% ============================================================

\begin{abstract}
\noindent
Este trabajo estudia el Aprendizaje por Refuerzo (RL) de forma progresiva, partiendo del paradigma más simple de toma de decisiones secuencial hasta los métodos de aproximación profunda con redes neuronales, a través de tres bloques de contenido interconectados.

\medskip
\noindent
\textbf{Capítulo 1 — Problema del Bandido de $k$-brazos.}
Se introduce el dilema fundamental del RL: el balance entre \textit{exploración} y \textit{explotación}. Un agente interactúa con un entorno estocástico de $k$ palancas, cada una con una distribución de recompensa desconocida, buscando maximizar la ganancia acumulada en un horizonte finito de $T = 300$ pasos. Se implementan y comparan empíricamente tres familias algorítmicas bajo $N = 500$ ejecuciones y distribuciones de recompensa Normal, Binomial y Bernoulli: (i)~\textbf{$\epsilon$-Greedy} y su variante con decaimiento $\epsilon$-Decay, que combinan explotación determinista con exploración aleatoria; (ii)~\textbf{UCB} (UCB1, UCB2, UCB1-Tuned), que racionalizan la exploración mediante límites de confianza superior sobre la incertidumbre de cada brazo; y (iii)~\textbf{Ascenso de Gradiente} (Softmax y Preferencias), que aprenden una distribución de selección sobre las acciones. Los resultados muestran que no existe un algoritmo universalmente óptimo: UCB2 sobresale en entornos Normales, mientras que el Gradiente de Preferencias ($\alpha = 0.5$) domina en el régimen binario Bernoulli con un arrepentimiento acumulado prácticamente nulo ($\text{Regret} \approx 0$).

\medskip
\noindent
\textbf{Capítulo 2 — Entornos Complejos: Métodos Tabulares.}
Se formula el problema de control óptimo como un Proceso de Decisión de Markov (MDP) $\langle \mathcal{S}, \mathcal{A}, \mathcal{P}, \mathcal{R}, \gamma \rangle$ y se evalúa la familia de métodos tabulares model-free sobre el entorno \texttt{CliffWalking-v0}. Se implementan Monte Carlo On/Off-Policy, SARSA, Expected SARSA, Q-Learning, Double Q-Learning y variantes de $n$-pasos. Se diseña una arquitectura modular que desacopla el agente, el algoritmo (\textit{Learner}) y el entorno (Gymnasium), garantizando reproducibilidad con $N = 10$ ejecuciones de 400 episodios ($\gamma = 0.95$, $\alpha = 0.5$). Los experimentos evidencian la superioridad de Expected SARSA ($\approx{-}15$ puntos de penalización) frente a Q-Learning ($\approx{-}56$) gracias a su robustez ante tasas de aprendizaje elevadas, y demuestran la ineficacia práctica de Monte Carlo en entornos con penalizaciones severas. Double Q-Learning resuelve el \textit{Maximization Bias} a costa de una convergencia más lenta. El análisis concluye con el límite intrínseco de los métodos tabulares: la \textit{maldición de la dimensionalidad} hace inviable la representación matricial $\mathbf{Q} \in \mathbb{R}^{|\mathcal{S}| \times |\mathcal{A}|}$ en espacios de estados continuos.

\medskip
\noindent
\textbf{Capítulo 3 — Entornos Complejos: Métodos Aproximados.}
Para superar la barrera de escalabilidad, se substituye la tabla $\mathbf{Q}$ por una función paramétrica $\hat{q}(\cdot\,;\,\boldsymbol{\theta})$ y se evalúan tres aproximadores de complejidad creciente sobre \texttt{CartPole-v1} ($\mathcal{S} \subseteq \mathbb{R}^4$) con $N = 10$ ejecuciones de 400 episodios: (i)~\textbf{SARSA Semi-Gradiente} con Tile Coding ($d = 16\,384$ características binarias, 32\,768 parámetros), cuya aproximación lineal limita la expresividad; (ii)~\textbf{DQN} (red $4 \to 64 \to 2$ con ReLU, Adam $\eta = 10^{-3}$, Experience Replay y red target), que combina la no linealidad neuronal con estabilización del proceso de Bellman; y (iii)~\textbf{Double DQN}, extensión de DQN que desacopla selección y evaluación de la acción óptima para corregir el sesgo de sobreestimación. Los resultados ($\overline{R}_{[-50:]}$: SARSA~69, DQN~218, Double DQN~189) confirman que la capacidad expresiva del aproximador determina la calidad de la política, que el Experience Replay es condición crítica de estabilidad, y que la ventaja de Double DQN sobre DQN solo emerge en horizontes de entrenamiento superiores a $T_\text{ep} > 1\,000$ episodios.

\medskip
\noindent
En conjunto, los tres capítulos trazan una progresión coherente desde la exploración-explotación en un único estado hasta el control de sistemas dinámicos continuos mediante redes neuronales profundas, ilustrando los compromisos fundamentales del RL entre \textbf{capacidad expresiva}, \textbf{eficiencia muestral} y \textbf{estabilidad del aprendizaje}.
\end{abstract}


% Índice
\tableofcontents

% Trabajos
\section{\texorpdfstring{Capítulo 1: Problema del Bandido de
\emph{k}-brazos}{Capítulo 1: Problema del Bandido de k-brazos}}\label{capuxedtulo-1-problema-del-bandido-de-k-brazos}

\subsection{1.1. Introducción}\label{introducciuxf3n}

\subsubsection{Descripción del Problema y
Relevancia}\label{descripciuxf3n-del-problema-y-relevancia}

El problema del bandido de \(k\)-brazos (o \emph{Multi-Armed Bandit})
ilustra formalmente el dilema fundamental del Aprendizaje por Refuerzo:
el balance iterativo entre la exploración y la explotación. Un agente
computacional se enfrenta a un entorno estacionario y estocástico
modelado como una máquina tragaperras con \(k\) palancas (brazos), cada
una ocultando una distribución de probabilidad de recompensa subyacente
diferente y desconocida \emph{a priori}.

Con un tiempo límite o presupuesto de \(T\) interacciones, la meta del
agente es maximizar su recompensa acumulativa. La toma de su decisión es
crítica: si confía únicamente en la palanca que le ha concedido mayor
ganancia económica hasta ahora (explotación), recae en el riesgo de
converger hacia una recompensa local ignorando una palanca numéricamente
superior. Si, por el contrario, invierte un esfuerzo desmedido en probar
nuevas palancas aleatorias para descubrir el verdadero margen de
beneficio (exploración), desperdicia iteraciones cobrando recompensas
bajas con la consecuente penalización temporal.

\subsubsection{Motivación del Trabajo}\label{motivaciuxf3n-del-trabajo}

La razón principal por la que estudiamos este problema es porque se
parece a muchas decisiones prácticas que tomamos en el mundo real (como
elegir en qué anuncio invertir dinero o qué tratamiento médico probar).
Es un entorno ``simple'' porque las decisiones que tomamos ahora no
cambian las opciones que tendremos en el futuro (cada tirada es
independiente). Sin embargo, estudiarlo es clave para entender cómo
funciona la inteligencia artificial cuando aprende a base de prueba y
error.

Si logramos comprender cómo estos algoritmos básicos
(\(\epsilon\)-Greedy o UCB) interactúan con la duda de \emph{``¿me quedo
con lo que ya sé que funciona o arriesgo a probar algo nuevo?''},
habremos construido los pilares necesarios para entender a los
verdaderos agentes robóticos complejos que resolverán problemas mucho
más difíciles en el futuro.

\subsubsection{Objetivos del Informe}\label{objetivos-del-informe}

El desarrollo de esta primera fase trata de cumplir los siguientes
objetivos técnicos: 1. Diseñar y programar un entorno modular en Python
capaz de modelar un Bandido bajo tres tipologías de distribución de
recompensa continuas y discretas: Normal, Binomial y Bernoulli. 2.
Implementar desde cero tres familias clásicas de algoritmos de toma de
decisiones: Explotación estocástica (\(\epsilon\)-Greedy y sus variantes
con decaimiento), Ascenso de Gradiente (Softmax y métricas de
Preferencia) y Optmismo basado en incertidumbre (UCB-1, UCB-2 y
UCB1-Tuned). 3. Evaluar dichas familias sometiéndolas empíricamente
frente al entorno mediante \(N=500\) ejecuciones independientes
simulando \(T=300\) pasos temporales. 4. Analizar la escalabilidad, la
ratio de acierto de selección óptima y la mitigación del arrepentimiento
acumulado (\emph{Cumulative Regret}).

\subsubsection{Organización del
Documento}\label{organizaciuxf3n-del-documento}

Para facilitar su asimilación, este primer capítulo estructura su
narrativa en el siguiente formato secuencial. En la sección de
\textbf{Desarrollo} (1.2), se sentarán las bases teóricas definiendo
matricial y probabilísticamente las distribuciones a las que se enfrenta
el bandido, adjuntando a la par un breve estado del arte de las
influencias bibliográficas. Seguidamente, en la sección de
\textbf{Algoritmos} (1.3), se expondrá la matemática formal y el
pseudocódigo subyacente de cada política heurística testeada. Las
asunciones y las simulaciones se materializan descriptivamente de la
mano de gráficas de rendimiento en la sección referente a la
\textbf{Evaluación/Experimentos} (1.4). Por último, las métricas son
discutidas críticamente en las \textbf{Conclusiones} (1.5) para
certificar el descubrimiento de los parámetros ganadores globales.

\subsection{1.2. Desarrollo}\label{desarrollo}

\subsubsection{Contexto y Antecedentes}\label{contexto-y-antecedentes}

El dilema de exploración frente a explotación es una encrucijadas
elemental dentro del Aprendizaje por Refuerzo. Mientras que en el
aprendizaje supervisado un algoritmo recibe explícitamente la respuesta
correcta a modo de etiqueta, en el aprendizaje evaluativo
(característico del RL) un agente debe descubrir autónomamente qué
acciones brindan la mayor ganancia probándolas. El modelo del
\emph{Multi-Armed Bandit} (Bandido Multibrazo) es el paradigma para
aislar y estudiar estas propiedades evaluativas y de decisión en un
entorno estacionario y no asociativo (es decir, donde el entorno no
cambia de estado). Comprender la balanza entre arriesgar ganancias
seguras presentes y el potencial de ganancias futuras desconocidas
cimienta la matemática que los agentes complejos aplican cuando operan
con los valores \(Q\) y variables de estado multidimensionales.

\subsubsection{Definición Formal del
Problema}\label{definiciuxf3n-formal-del-problema}

El entorno estocástico se define bajo el espacio de recompensas
discretas o continuas para un conjunto finito de acciones o brazos
\(a \in \mathcal{A} = \{1, 2, \dots, k\}\).

El verdadero y desconocido ``valor esperado'' de elegir una acción \(a\)
se denota habitualmente como el valor a maximizar \(q_*(a)\), y está
formulado como la esperanza matemática de la recompensa en un instante
\(t\): \[ q_*(a) = \mathbb{E}[R_t \mid A_t = a] \]

A lo largo del presente marco de validación, las recompensas \(R_t\)
derivan su valor estocástico basándose en tres entornos implementados:

\begin{enumerate}
\def\labelenumi{\arabic{enumi}.}
\item
  \textbf{Entorno de Recompensas Normales (Gausianas)} La retribución
  estocástica de cada acción se muestrea de una distribución diferencial
  \(R_t \sim \mathcal{N}(\mu_a, \sigma_a^2)\), donde la verdadera medida
  subyacente de recompensa está dominada por su función de densidad de
  probabilidad (PDF):
  \[ f(x \mid \mu_a, \sigma^2) = \frac{1}{\sqrt{2\pi\sigma^2}} e^{-\frac{(x - \mu_a)^2}{2\sigma^2}} \]
\item
  \textbf{Entorno de Recompensas Binomiales} En este escenario, la
  activación del brazo \(a\) devuelve una recompensa entera discreta
  (\(x \in \mathbb{N}\)) que contabiliza la cantidad total de éxitos
  sobre \(n_a\) ensayos independientes. Es dirigido por la variable de
  masa probabilística \(R_t \sim B(n_a, p_a)\):
  \[ P(R_t=x) = \binom{n_a}{x} p_a^x (1 - p_a)^{n_a - x} \]
\item
  \textbf{Entorno de Recompensas de Bernoulli} Un escenario más sencillo
  donde el resultado de tirar de la palanca solo puede ser un éxito o un
  fracaso absoluto (\(R_t \in \{0, 1\}\)). Modela la variable
  \(R_t \sim Bernoulli(p_a)\) en donde el valor promedio de iteración a
  estabilizar es exactamente la probabilidad intrínseca al brazo
  evaluado (\(q_*(a) = p_a\)):
  \[ P(R_t) = p_a^{R_t} (1 - p_a)^{1 - R_t} \]
\end{enumerate}

\subsubsection{Enfoque y Justificación
Metodológica}\label{enfoque-y-justificaciuxf3n-metodoluxf3gica}

Para abordar el problema estocástico, se ha abstraído un
\emph{framework} orientado a objetos que desacopla completamente el
``Entorno'' del ``Agente''. Esto permite insertar módulos de políticas
que basan su decisión algorítmica interactuando con las mismas
estructuras pre-configuradas \(R_t\). El presente trabajo se caracteriza
por implementar la comparativa simultánea y simétrica del estado del
arte (\emph{\(\epsilon\)-Greedy} frente al Límite de Confianza
\emph{UCB} y el cálculo heurístico \emph{Softmax}) bajo un volumen
estadístico (\(N=500\) iteraciones promediadas), garantizando un estudio
numérico de optimización con mitigación del \emph{Arrepentimiento
Acumulado} en ecosistemas tanto lógicos continuos como discretos.

\subsubsection{1.2.1. Trabajos
Relacionados}\label{trabajos-relacionados}

El planteamiento algorítmico, su formulación paramétrica y las
heurísticas estandarizadas analizadas en este proyecto se fundamentan y
se derivan de la literatura existente en el campo del Aprendizaje por
Refuerzo: - \textbf{Métodos Epsilon y Ascenso de Gradiente (Sutton y
Barto, 2018):} El diseño de las estrategias \(\epsilon\)-Greedy y de los
métodos probabilísticos como el \emph{Ascenso de Gradiente} se ha
construido siguiendo las reglas y fundamentos explicados en el libro de
referencia ``Reinforcement Learning: An Introduction''
\cite{sutton2018reinforcement}. - \textbf{Manejo de la Incertidumbre con
UCB (Auer et al., 2002):} Para programar los algoritmos UCB (Upper
Confidence Bound) de forma robusta y entender cómo utilizan sus fórmulas
matemáticas para explorar las opciones menos conocidas, nos hemos
apoyado en el estudio ``Finite-time Analysis of the Multiarmed Bandit
Problem'' \cite{auer2002finitetime}.

\subsection{1.3. Algoritmos}\label{algoritmos}

Para resolver el problema del bandido de \(k\)-brazos hemos implementado
tres enfoques distintos: métodos que combinan explotar lo conocido con
explorar al azar, algoritmos que premian probar opciones inciertas para
dejar de dudar, y tácticas que eligen las acciones basándose en
preferencias relativas en lugar de en ganancias directas. A continuación
se justifican sus enfoques y se detalla el funcionamiento lógico de los
pilares principales de cada familia algorítmica estudiada.

\subsubsection{\texorpdfstring{Familia 1: Lógicas de
Exploración-Explotación
(\(\epsilon\)-Greedy)}{Familia 1: Lógicas de Exploración-Explotación (\textbackslash epsilon-Greedy)}}\label{familia-1-luxf3gicas-de-exploraciuxf3n-explotaciuxf3n-epsilon-greedy}

\textbf{Justificación:} El algoritmo \(\epsilon\)-Greedy es el punto de
partida más clásico en el Aprendizaje por Refuerzo. Su funcionamiento es
muy sencillo y rápido de calcular: la mayor parte del tiempo (con
probabilidad \(1-\epsilon\)) el agente elige la opción que hasta ahora
le ha dado más beneficios, pero de vez en cuando (con una pequeña
probabilidad \(\epsilon\)) elige una opción al azar para explorar. A
partir de aquí nace una mejora llamada \(\epsilon\)-Decay, que soluciona
el problema de explorar para siempre: va reduciendo poco a poco la
probabilidad de explorar (multiplicándola por un factor \(\lambda\)), de
modo que, a medida que el agente aprende, se vuelve más seguro y se
centra en asegurar las ganancias.

\textbf{Pseudocódigo de \(\epsilon\)-Greedy Clásico:}

\begin{algorithm}
\caption{Algoritmo $\epsilon$-Greedy para K-Brazos}
\begin{algorithmic}[1]
\REQUIRE $\epsilon \in [0, 1]$, Total de acciones $K$
\STATE Inicializar $Q(a) \leftarrow 0$ para todo $a=1,\dots, K$
\STATE Inicializar $N(a) \leftarrow 0$ para todo $a=1,\dots, K$
\FOR{cada paso $t=1, 2, \dots, T$}
    \STATE $p \leftarrow$ Valor aleatorio de Distribución Uniforme(0, 1)
    \IF{$p < \epsilon$}
        \STATE $A_t \leftarrow$ Acción aleatoria uniforme en $\{1, \dots, K\}$
    \ELSE
        \STATE $A_t \leftarrow \arg\max_a Q(a)$  \COMMENT{Ruptura de empates al azar}
    \ENDIF
    \STATE Ejecutar $A_t$, observar recomensa $R_t$
    \STATE $N(A_t) \leftarrow N(A_t) + 1$
    \STATE $Q(A_t) \leftarrow Q(A_t) + \frac{1}{N(A_t)} \big[R_t - Q(A_t)\big]$
\ENDFOR
\end{algorithmic}
\end{algorithm}

\subsubsection{Familia 2: Optimismo ante la Incertidumbre
(UCB)}\label{familia-2-optimismo-ante-la-incertidumbre-ucb}

\textbf{Justificación:} A diferencia de explorar siempre al azar, los
métodos de \emph{Upper Confidence Bound} (UCB) son curiosos de forma
mucho más inteligente. UCB calcula cuánta confianza tiene en lo que ya
sabe de cada palanca. Si una opción se ha probado muy poco, su nivel de
incertidumbre es alto, así que el algoritmo decide elegirla para ``salir
de dudas''. Esto convierte a UCB en un método robusto y bueno para no
perder el tiempo. En este trabajo hemos estudiado el algoritmo básico
UCB1 (que usa un número \(c\) para controlar de forma fija esa
curiosidad), el modelo UCB2 (que optimiza el proceso probando las
opciones en ``ráfagas'' o bloques cada vez más largos en lugar de
recalcular a cada paso), y la variante avanzada UCB1-Tuned (que se
ajusta a sí misma midiendo la variabilidad de los premios que recibe).

\textbf{Pseudocódigo de UCB2:}

\begin{algorithm}
\caption{Algoritmo UCB2}
\begin{algorithmic}[1]
\REQUIRE Constante exploratoria $\alpha \in (0, 1)$
\STATE Inicializar contador de épocas por brazo: para cada $a$, $r_a \leftarrow 0$
\STATE Iniciar explorando todos los brazos una vez, actualizando su recompensa promedio $Q(a)$
\FOR{paso de tiempo $t = K+1, \dots$}
    \STATE Seleccionar brazo $A_t \leftarrow \arg\max_a \left[ Q(a) + \sqrt{\frac{(1 + \alpha) \ln (e \cdot t / \tau(r_a))}{2 \tau(r_a)}} \right]$ \newline
           (donde $\tau(r) = \lceil (1+\alpha)^r \rceil$)
    \STATE Determinar longitud de la ráfaga: $\Delta \leftarrow \tau(r_{A_t} + 1) - \tau(r_{A_t})$
    \FOR{iteraciones $i = 1$ hasta $\Delta$}
        \STATE Ejecutar $A_t$, percibir recompensa $R$
        \STATE $N(A_t) \leftarrow N(A_t) + 1$
        \STATE $Q(A_t) \leftarrow Q(A_t) + \frac{1}{N(A_t)} \big[R - Q(A_t)\big]$
    \ENDFOR
    \STATE Actualizar época del brazo: $r_{A_t} \leftarrow r_{A_t} + 1$
\ENDFOR
\end{algorithmic}
\end{algorithm}

\subsubsection{Familia 3: Ascenso de Gradiente Estocástico (Softmax y
Preferencias)}\label{familia-3-ascenso-de-gradiente-estocuxe1stico-softmax-y-preferencias}

\textbf{Justificación:} Si los algoritmos solo se fijan en los puntos
exactos de cada opción (\(Q\)), pueden tener problemas para decidir si
los premios son muy parecidos (por ejemplo, ganar \(+85\) o \(+84\)).
Los métodos basados en el Gradiente lo resuelven usando ``Preferencias''
(\(H_a\)). En vez de memorizar valores numéricos rígidos, simplemente
aprenden qué opción prefieren en relación al resto y deciden usando
porcentajes (con la Ley Exponencial Suave o Gibbs). En nuestras pruebas,
esta estrategia del \emph{Gradiente de Preferencias} trataremos de
demostrar si es la que mejor se adapta a los entornos donde las
recompensas son simplemente ``todo o nada'' (como sucede en el modelo de
\emph{Bernoulli}).

\textbf{Pseudocódigo de Base Gibbs / Gradiente:}

\begin{algorithm}
\caption{Algoritmo de Variación de Gradiente de Preferencias}
\begin{algorithmic}[1]
\REQUIRE Tasa de aprendizaje $\alpha > 0$
\STATE Inicializar preferencias $H(a) \leftarrow 0 \;\; \forall a$ 
\STATE Inicializar recompensa promedio de validación $\bar{R} \leftarrow 0$
\FOR{$t = 1, \dots, T$}
    \STATE Calcular probabilidades de la curva Softmax:\\
    $P(a) \leftarrow \frac{\exp(H(a))}{\sum_{b=1}^K \exp(H(b))} \quad \forall a$
    \STATE Elegir acción $A_t$ en base a la distribución de densidad $P(\cdot)$
    \STATE Obtener recompensa $R_t$
    \STATE Actualizar la base empírica global $\bar{R} \leftarrow \bar{R} + \frac{1}{t}(R_t - \bar{R})$
    \FORALL{$a \in \{1, \dots, K\}$}
        \IF{$a = A_t$}
            \STATE $H(a) \leftarrow H(a) + \alpha(R_t - \bar{R})(1 - P(a))$
        \ELSE
            \STATE $H(a) \leftarrow H(a) - \alpha(R_t - \bar{R})P(a)$
        \ENDIF
    \ENDFOR
\ENDFOR
\end{algorithmic}
\end{algorithm}

\subsection{1.4. Evaluación y
Experimentos}\label{evaluaciuxf3n-y-experimentos}

\subsubsection{1.4.1. Configuración
Experimental}\label{configuraciuxf3n-experimental}

Para garantizar la validez científica y la total reproducibilidad de los
resultados a lo largo del estudio de las tres familias de distribuciones
(Normal, Binomial y Bernoulli), el entorno de pruebas fue estandarizado
con los siguientes hiperparámetros de simulación global: -
\textbf{Ejecuciones (\emph{Runs}):} \(N = 500\) pruebas independientes
estocásticas. Las curvas finales presentadas corresponden a la matriz
promediada probabilísticamente sobre este número total de semillas. -
\textbf{Horizonte de Simulador (\emph{Steps}):} \(T = 300\) secuencias
temporales de decisión por ejecución. - \textbf{Semilla Base Estática:}
Se inyectó globalmente una semilla matemática (aleatoria constante,
habitualmente \texttt{np.random.seed(1024)}) para inicializar la
topología generatriz del bandido de forma idéntica ante la comparativa
de los tres algoritmos.

Para llevar a cabo las simulaciones, se definieron los siguientes
espectros de hiperparámetros iniciales para cada modelo inteligente
(Tabla \ref{tab:Parametros}):

\begin{table*}[t]
\centering
\caption{Hiperparámetros definidos y comparados por familia algorítmica.}
\label{tab:Parametros}
\begin{tabular}{|l|l|l|}
\hline
\textbf{Familia Algorítmica} & \textbf{Configuración / Algoritmo} & \textbf{Hiperparámetros Evaluados} \\ \hline
Exploración-Explotación & $\epsilon$-Greedy & $\epsilon \in \{0.0, 0.01, 0.1\}$ \\ \hline
Exploración-Explotación & $\epsilon$-Decay & Tasa de decaimiento $\lambda \in \{0.8, 0.99, 0.9999\}$ \\ \hline
Ascenso de Gradiente & Softmax & Temperatura $\tau \in \{0.5, 1.0, 2.0\}$ \\ \hline
Ascenso de Gradiente & Gradiente de Preferencias & Tasa de aprendizaje $\alpha \in \{0.5, 1.0, 2.0\}$ \\ \hline
Límite Sup. Condianza & UCB1 & Constante de exploración $c \in \{0.5, 1.0, 1.5\}$ \\ \hline
Límite Sup. Condianza & UCB2 & Expansión de confianza temporal $\alpha \in \{0.25, 0.5, 0.75\}$ \\ \hline
Límite Sup. Condianza & UCB1-Tuned & Ninguno (Varianza de la Recompensa empírica) \\ \hline
\end{tabular}
\end{table*}

\subsubsection{1.4.2. Métricas de
Desempeño}\label{muxe9tricas-de-desempeuxf1o}

La evaluación del rendimiento integral de los agentes propuestos se
analizó multidimensionalmente observando su desempeño iterativo bajo 4
prismas métricos: 1. \textbf{Recompensa Promedio Continua (\emph{Average
Reward}):} Traza iterativa de la estimación \(\bar{R_t}\); refleja la
capacidad del agente de percibir y adaptar ganancia bruta promedio. 2.
\textbf{Ratio Probabilístico de Selección (\emph{\% Optimal Action}):}
Discrimina de forma binaria el acierto o el fallo explícito al localizar
matemáticamente la mejor máquina subyacente \(\arg\max q_*(a)\). 3.
\textbf{Arrepentimiento Acumulado (\emph{Cumulative Regret}):} Mide la
acumulación de pérdida teórica u ``oportunidad perdida'' a nivel
temporal de decantarse por opciones penalizadoras frente a un escenario
``clarividente''. Se busca priorizar gráficas cóncavas y rápidamente
estabilizadas (planas). 4. \textbf{Histogramas Distribucionales:}
Estadísticos puros del muestreo frecuentista con el comportamiento
latente final del ecosistema (\emph{Arm Statistics}).

\subsubsection{1.4.3. Análisis de Resultados (Gráficas y
Comparativas)}\label{anuxe1lisis-de-resultados-gruxe1ficas-y-comparativas}

La convergencia de rendimiento obtenida tras los experimentos numéricos
extraídos de \texttt{epsilon\_greedy.ipynb},
\texttt{ascenso\_gradiente.ipynb} y \texttt{UCB\_algorithms.ipynb} ha
posibilitado la categorización de parámetros óptimos para las tres ramas
del estado del arte, englobados en la siguiente matriz global (Tabla
\ref{tab:Resultados}):

\begin{table*}[t]
\centering
\caption{Parámetros con mayor rendimiento cruzados por distribución.}
\label{tab:Resultados}
\begin{tabular}{|l|l|p{7cm}|}
\hline
\textbf{Familia} & \textbf{Mejor Hiperparámetro} & \textbf{Análisis Crítico} \\ \hline
$\epsilon$-Greedy   & $\epsilon = 0.1$                                  & Globalmente superior a versiones estáticas conservadoras ($\epsilon=0.01$). \\ \hline
Decaimento      & $\lambda = 0.99$                                  & Combina convergencia exponencial de exploración. Cae veloz frente a entornos como Bernoulli si su suelo ($\epsilon_{min}$) es severo.\\ \hline
Softmax         & Temperatura $\tau \in [0.5, 1.0]$                 & Curva modesta en continuas; catastrófico (20\% acierto) en dominios Bernoulli. \\ \hline
P. Gradiente    & Ratio de Aprendizaje $\alpha \in [0.5, 1.0]$      & La revelación. Su matriz de distancias absolutas mitiga por completo entornos acotados y discretos (Bernoulli).  \\ \hline
UCB             & $c = 1.0$ (o $c=0.5$ Bernoulli)                    & Curvas de convergencia teórica muy seguras y logarítmicas. Rinde maravillosamente para incertidumbres Gausianas. \\ \hline
UCB1-Tuned      & Ninguno (No paramétrico)                          & Defensa inquebrantable en todas las distribuciones; el entorno escala su nivel de exploración por la varianza natamente. \\ \hline
\end{tabular}
\end{table*}

\subsubsection{\texorpdfstring{Resultados Familia
\(\epsilon\)-Greedy}{Resultados Familia \textbackslash epsilon-Greedy}}\label{resultados-familia-epsilon-greedy}

\begin{itemize}

\item
  \textbf{Decaimiento (\(\epsilon\)-Decay):} Aunque en teoría debería
  ser mejor, en la práctica el decaimiento no ha supuesto ninguna
  ventaja clara en los entornos \emph{Normal} y \emph{Binomial}. De
  hecho, ha funcionado igual o ligeramente \textbf{peor} que el modelo
  básico (\(\epsilon\)-Greedy puro). De todos los ajustes probados,
  reducir la probabilidad de explorar progresivamente con un
  multiplicador \(\lambda=0.99\) fue lo más equilibrado. En cambio, si
  la exploración se reduce muy despacio (\(\lambda=0.9999\)), el
  algoritmo pierde demasiado tiempo escogiendo opciones al azar,
  acumulando muchas pérdidas.
\item
  \textbf{Bernoulli}: En este caso particular, hemos experimentado un
  problema grave: si frenamos la caída de exploración demasiado pronto
  (fijando un mínimo inamovible de \(\epsilon_{min}=0.001\)), el
  algoritmo deja de aprender de golpe. Esto provoca que se estanque
  prematuramente y no logre igualar la enorme tasa de aciertos que sí
  consiguen otros competidores.
\end{itemize}

\subsubsection{Resultados Familia de Ascenso de
Gradiente}\label{resultados-familia-de-ascenso-de-gradiente}

\begin{itemize}

\item
  \textbf{Softmax frente a los entornos:} Mientras que el método Softmax
  se defiende de forma aceptable en entornos donde los premios varían de
  forma continua (distribución \emph{Normal}), fracasa completamente
  cuando los premios son de ``todo o nada'' (\emph{Bernoulli}). En este
  último caso, el algoritmo se atasca rápidamente, eligiendo la mejor
  opción apenas un 20\% de las veces y acumulando muchísimas pérdidas a
  lo largo del tiempo.
\item
  \textbf{Gradiente de Preferencias vs Softmax:} Este ha sido uno de los
  grandes descubrimientos de las pruebas. El \emph{Gradiente de
  Preferencias}, al guiarse por ``cuánto le gusta'' una opción
  (porcentajes de preferencia) en lugar de intentar memorizar valores
  exactos de puntos, \textbf{ha superado con mucha claridad a Softmax en
  todos los escenarios}. Su éxito es especialmente espectacular en el
  entorno \emph{Bernoulli}, donde consigue elegir la mejor palanca casi
  el 100\% de las veces desde muy temprano. Esto hace que sus pérdidas
  (arrepentimiento) se queden en unos ridículos 3 puntos tras 300
  intentos (frente a los 70 puntos perdidos por Softmax).
\end{itemize}

\subsubsection{Resultados Familia UCB}\label{resultados-familia-ucb}

\begin{itemize}

\item
  \textbf{UCB1 vs.~UCB2:} El modelo clásico UCB1 aprende bastante bien
  en casi todas las situaciones probadas, sobre todo cuando su nivel de
  confianza inicial está fijado en \(c=1.0\) (para los modelos
  \emph{Normal} y \emph{Binomial}). Sin embargo, en el juego de premios
  binarios (\emph{Bernoulli}) le va mejor si arranca más precavido con
  un \(c=0.5\). Su hermano mayor, \textbf{UCB2}, muestra dos caras
  opuestas: \textbf{falla estrepitosamente en los entornos de Binomial y
  Bernoulli}, donde se atasca eligiendo constantemente la peor opción
  sin aprender casi nada, pero a cambio resulta ser rapidísimo y el
  mejor de todos cuando los premios son más naturales y variados
  (entornos \emph{Normales}).
\item
  \textbf{UCB1-Tuned}: Esta versión con ``piloto automático'' consiguió
  igualar algunos de los buenos resultados de UCB1 sin necesidad de
  configurar ningún parámetro a mano. Aún así, sigue pasándolo bastante
  mal cuando las reglas del juego son drásticas (como en el entorno
  \emph{Bernoulli}), donde necesita atascarse al menos 100 intentos
  seguidos antes de lograr darse cuenta de cuál es el brazo ganador.
\end{itemize}

Para ver todo esto de forma mucho más clara, las siguientes gráficas
(sacadas de probar los algoritmos en entornos de premios de tipo
\emph{Normal}) muestran un resumen visual perfecto. Tal y como
descubrimos al analizar a fondo las simulaciones, \textbf{la gráfica del
Arrepentimiento Acumulado resulta ser siempre la más útil y fácil de
entender}. Esto se debe a que, a diferencia de otras medidas, en el
Arrepentimiento se puede distinguir a simple vista si un algoritmo sigue
aprendiendo y reduciendo sus pérdidas o si ya se ha estancado por
completo.

\begin{figure}[h]
    \centering
    \includegraphics[width=\columnwidth]{k_brazos/img/eg_regret_normal.png}
    \caption{Curva de Arrepentimiento Acumulado para la familia $\epsilon$-Greedy.}
    \label{fig:egreedy_perf}
\end{figure}

En la Figura \ref{fig:egreedy_perf}, se discierne claramente cómo una
caída suave del decaimiento temporal (\(\lambda = 0.99\)) permite
achatar asintóticamente el \emph{regret} mitigando estancamientos
lineales que sí sufrían ramas no decadentes.

\begin{figure}[h]
    \centering
    \includegraphics[width=\columnwidth]{k_brazos/img/ag_regret_normal.png}
    \caption{Curva de Arrepentimiento Acumulado comparando Preferencias vs Softmax.}
    \label{fig:ag_perf}
\end{figure}

Por otro lado, la comparación que vemos en la Figura \ref{fig:ag_perf}
nos demuestra algo muy interesante: aunque el algoritmo Softmax arranca
bien y consigue ganancias rápidamente, basar las decisiones puramente en
las \textbf{Preferencias} de cada opción frente al resto (la curva
verde) resulta ser una estrategia absolutamente imbatible a medio y
largo plazo siempre que se tenga configurada la tasa de aprendizaje
correcta (\(\alpha=1\)).

\begin{figure}[h]
    \centering
    \includegraphics[width=\columnwidth]{k_brazos/img/ucb_regret_normal.png}
    \caption{Curva de Arrepentimiento Acumulado demostrando el inquebrantable éxito de UCB1.}
    \label{fig:ucb_perf}
\end{figure}

Por último, la Figura \ref{fig:ucb_perf} demuestra en la práctica lo
bien que funciona la idea principal del método UCB: si castigamos
fuertemente a las opciones que ya hemos probado demasiadas veces,
conseguimos que el algoritmo frene en seco sus pérdidas
(arrepentimiento). Todo esto lo logra con una precisión asombrosa, de
forma rápida y sin necesidad de fórmulas ni ajustes complicados, tal y
como atestigua el excelente recorrido de la variante con ``piloto
automático'' (\emph{UCB1-Tuned} en la curva verde).

\subsection{1.5. Conclusiones}\label{conclusiones}

\subsubsection{Resumen de Principales
Resultados}\label{resumen-de-principales-resultados}

Tras probar a fondo todos los algoritmos en el problema del bandido de
k-brazos, sacamos una conclusión clara: no existe un algoritmo perfecto
para todo. La mejor estrategia siempre dependerá del tipo de entorno y
de cómo se comporten los premios.

\begin{itemize}

\item
  En entornos donde las recompensas varían de forma continua y lógica
  (como en la distribución \emph{Normal}), métodos como \textbf{UCB1} y
  especialmente \textbf{UCB2} aprenden rapidísimo al elegir de forma
  inteligente qué opciones explorar. Sin embargo, \textbf{solo UCB1}
  logra mantener ese buen nivel en entornos algo más rígidos como el
  \emph{Binomial}, ya que UCB2 se atasca por completo.
\item
  Cuando nos enfrentamos a escenarios drásticos de ``todo o nada''
  (recompensas \emph{Bernoulli}), casi todos los algoritmos anteriores
  fallan de forma estrepitosa y acumulan pérdidas. En estos casos tan
  extremos, el ganador es el \textbf{Gradiente de Preferencias}. En vez
  de obsesionarse con calcular exactamente cuántos puntos da cada
  opción, este modelo se limita a comparar qué tan buena es una palanca
  respecto de las demás; gracias a esto, logra dominar el juego y no
  equivocarse casi nunca (\(Regret \approx 0\)).
\item
  Para situaciones variadas o cuando no tenemos mucha potencia de
  cálculo, recurrir a los clásicos métodos semialeatorios (como
  \textbf{\(\epsilon\)-Decay} o directamente un \(\epsilon\) fijo de
  \(0.1\)) resulta ser una estrategia buena.
\end{itemize}

\subsubsection{Limitaciones del Estudio}\label{limitaciones-del-estudio}

La limitación principal de los resultados reportados en este documento
es que el juego asume que las reglas no cambian, la propiedad de
\textbf{estacionalidad} del modelo K-Brazos abstraído. En el mundo real,
la característica estocástica y las medias de recompensa (\(\mu\)) de
una máquina o variable evolucionan y mutan con el paso del tiempo. Si
transladamos estos algoritmos a la vida real, estos algoritmos podrían
fracasar estrepitosamente porque se obcecarían en descartar opciones de
escenarios que de repente podrían volverse muy rentables.

Además, este marco experimentacional excluye variables de ``contexto''
(variables climáticas, visuales o previas a la decisión) de cada tirada
o instante \(t\), tratando a las palancas bajo el paradigma
probabilístico absoluto e idéntico en cada \emph{step}.

\subsubsection{Impacto y Reflexión en el Aprendizaje por
Refuerzo}\label{impacto-y-reflexiuxf3n-en-el-aprendizaje-por-refuerzo}

Estudiar a fondo estas tres familias de algoritmos del bandido (apostar
a lo seguro, investigar las opciones menos conocidas, o guiarse por las
preferencias) no es para nada un ejercicio anticuado o puramente
teórico. Al contrario: el concepto del \emph{Arrepentimiento} y cómo
minimizarlo son los cimientos más básicos que le dicen a cualquier
Inteligencia Artificial cómo debe tomar sus decisiones.

Estas fórmulas e ideas matemáticas conforman el motor interno inicial de
decisiones de cualquier agente moderno y complejo en el mundo real (como
los sistemas de \emph{Deep Reinforcement Learning} de las IA generativas
o de los coches autónomos).

\subsubsection{Líneas Futuras de
Estudio}\label{luxedneas-futuras-de-estudio}

El siguiente paso natural para avanzar en este campo sería crear
simulaciones de \textbf{Bandidos No Estacionarios (\emph{Non-stationary
Bandits})}, es decir, enseñar a la Inteligencia Artificial a jugar en
entornos donde las reglas y los premios van cambiando con el tiempo.
Para lograrlo, habría que dotar a los algoritmos con capacidades de
adaptación ``olvidadizas'': mecanismos matemáticos que les permitan
ignorar lecciones antiguas que ya no sirven y centrarse en lo que está
ocurriendo ahora mismo.

Más adelante, el objetivo final sería evolucionar hacia \textbf{Bandidos
Contextuales (\emph{Contextual Bandits})}. En este escenario, la IA no
jugaría a ciegas, sino que recibiría ``pistas'' del entorno antes de
tener que elegir una opción (por ejemplo, saber qué tiempo hace antes de
recomendar un tipo de ropa). Esto abre directamente las puertas hacia
sistemas de inteligencia artificial mucho más completos y complejos que
los que hemos visto hasta ahora.

\input{metodos_tabulares}
% ============================================================
%  Capítulo 3 — Aprendizaje en Entornos Complejos: Métodos Aproximados
%  EML-RL · Lucas Ortiz ·
% ============================================================

\chapter{Entornos Complejos: Métodos Aproximados}

% ============================================================
\section{Introducción}
% ============================================================

\subsection{Descripción del Problema y su Relevancia}

El \textbf{Aprendizaje por Refuerzo} (RL, del inglés \textit{Reinforcement Learning}) estudia cómo un agente autónomo aprende a tomar decisiones óptimas interactuando con un entorno mediante el esquema de prueba y error. En su formulación estándar como \textbf{Proceso de Decisión de Markov} (MDP), el agente observa el estado del entorno $s_t \in \mathcal{S}$, selecciona una acción $a_t \in \mathcal{A}$ y recibe una recompensa escalar $r_{t+1} \in \mathbb{R}$, con el objetivo de maximizar el retorno acumulado descontado:

\begin{equation}
  G_t = \sum_{k=0}^{\infty} \gamma^k R_{t+k+1}, \qquad \gamma \in [0, 1)
\end{equation}

Cuando el espacio de estados $\mathcal{S}$ es \textbf{continuo} o de alta dimensionalidad, la representación tabular de la función de valor-acción $Q: \mathcal{S} \times \mathcal{A} \to \mathbb{R}$ resulta computacionalmente inviable. El problema central de los \textit{entornos complejos} reside en cómo generalizar el conocimiento adquirido en estados visitados hacia estados no visitados de la misma región del espacio.

La relevancia de este problema es doble: por un lado, la mayoría de las aplicaciones industriales y científicas de RL involucran espacios de estados continuos (robótica, vehículos autónomos, gestión de energía, biomedicina); por otro, los avances en \textbf{aproximación de funciones} —especialmente mediante redes neuronales profundas— han impulsado logros como el control a nivel humano en videojuegos Atari \cite{Mnih2015} o el dominio de juegos de mesa complejos \cite{Silver2017}.

\subsection{Motivación del Trabajo}

El punto de partida de este capítulo es el límite de escalabilidad evidenciado en el Capítulo 2 (métodos tabulares). Dado un entorno como \texttt{MountainCar-v0} con un espacio de estados $\mathcal{S} \subseteq \mathbb{R}^2$, la discretización en una rejilla de $20 \times 20 = 400$ celdas muestra que el agente requiere un presupuesto de episodios muy elevado para mejorar de forma apreciable, ya que:

\begin{equation}
  \hat{Q}(s, a) \text{ solo actualiza la celda } \lfloor s / \Delta \rfloor, \quad \text{sin generalización a } s \pm \epsilon
\end{equation}

Esta rigidez estructural motiva el estudio de \textbf{métodos aproximados}, que sustituyen la tabla por una función paramétrica $\hat{q}(s, a\,;\, \boldsymbol{\theta})$ capaz de interpolar entre estados no visitados. La motivación específica de este trabajo es:

\begin{enumerate}
  \item \textbf{Superar la barrera de escalabilidad} de los métodos tabulares en entornos con observaciones continuas.
  \item \textbf{Comparar el espectro de aproximadores}: desde la aproximación lineal con Tile Coding hasta las redes neuronales profundas (DQN, Double DQN).
  \item \textbf{Caracterizar el compromiso} entre capacidad expresiva del aproximador, estabilidad del entrenamiento y eficiencia muestral.
  \item \textbf{Evaluar el papel de la infraestructura Gymnasium} como marco estándar de experimentación reproducible en RL.
\end{enumerate}

\subsection{Objetivos del Informe}

Este capítulo persigue los siguientes objetivos:

\begin{enumerate}
    \item \textbf{O1}. Formalizar el marco matemático de la aproximación de funciones en RL.
    \item \textbf{O2}. Describir e implementar tres familias de métodos aproximados: SARSA Semi-Gradiente con Tile Coding, DQN y Double DQN.
    \item \textbf{O3}. Diseñar y ejecutar un protocolo experimental reproducible sobre \texttt{CartPole-v1}.
    \item \textbf{O4}. Analizar cuantitativamente los resultados y comparar las familias algorítmicas.
    \item \textbf{O5}. Extraer conclusiones sobre la adecuación de cada aproximador.
\end{enumerate}

\subsection{Organización del Capítulo}

El resto del capítulo se estructura como sigue: \S3.2 presenta la definición formal del problema y los trabajos relacionados; \S3.3 detalla los pseudocódigos y ecuaciones de los tres algoritmos; \S3.4 describe la configuración experimental, métricas y resultados; \S3.5 sintetiza los hallazgos, limitaciones y líneas futuras.

% ============================================================
\section{Desarrollo}
% ============================================================

\subsection{Trabajos Relacionados}

La aproximación de funciones en RL tiene sus raíces en los trabajos fundacionales de Sutton (1988) sobre el aprendizaje por Diferencias Temporales (TD) y la tesis de Watkins (1989) sobre Q-Learning.

\subsubsection{Aproximación Lineal y Codificación de Características}

Sutton \& Barto \cite{Sutton2018} formalizan el \textbf{SARSA Semi-Gradiente} con aproximación lineal, demostrando convergencia al óptimo proyectado bajo hipótesis de función de distribución de visita estacionaria $\mu(s)$. El uso de \textbf{Tile Coding} \cite{Sutton1996} como esquema de representación permite aprovechar la sparsidad del vector de características: solo $m$ de los $d$ componentes son activos en cada paso, reduciendo la evaluación de $O(d)$ a $O(m)$.

\subsubsection{Deep Q-Network (DQN)}

Mnih et al.~\cite{Mnih2015} introdujeron \textbf{DQN}, que combina Q-Learning con redes neuronales convolucionales para aprender directamente desde píxeles en 49 juegos Atari. Sus contribuciones técnicas clave fueron:

\begin{itemize}
  \item \textbf{Experience Replay} \cite{Lin1992}: almacenamiento de transiciones en un buffer $\mathcal{D}$ y muestreo aleatorio de mini-batches para romper correlaciones temporales.
  \item \textbf{Red Target} (\textit{target network}): copia periódica de los pesos $\boldsymbol{\theta}^- \leftarrow \boldsymbol{\theta}$ para estabilizar los targets de Bellman durante el entrenamiento.
\end{itemize}

\subsubsection{Double DQN}

Van Hasselt et al.~\cite{vanHasselt2016} identificaron que DQN estándar sobreestima sistemáticamente los valores $Q$ debido al operador $\max$ en el cálculo del target:

\begin{equation}
  \mathbb{E}\!\left[\max_{a'} Q_{\boldsymbol{\theta}}(s', a')\right] \geq \max_{a'} \mathbb{E}\!\left[Q_{\boldsymbol{\theta}}(s', a')\right]
\end{equation}

Su solución, \textbf{Double DQN}, desacopla la \textit{selección} de acción (red online $\boldsymbol{\theta}$) de la \textit{evaluación} de su valor (red target $\boldsymbol{\theta}^-$).

\subsubsection{Otras Referencias Relevantes}

Véase la Tabla~\ref{tab:extensiones}.

\begin{table*}[htb!]
\caption{Principales referencias de DQN en la literatura}\label{tab:extensiones}
\centering\small
\begin{tabular}{@{}llp{8cm}@{}}
\toprule
\textbf{Método} & \textbf{Referencia} & \textbf{Contribución principal} \\
\midrule
Dueling DQN & Wang et al.~(2016) & Desacopla $V(s)$ y $A(s,a)$ \\
Prioritized Replay & Schaul et al.~(2016) & Muestreo proporcional al error TD \\
A3C & Mnih et al.~(2016) & Actor-Critic asíncrono sin replay buffer \\
Rainbow & Hessel et al.~(2018) & Combinación de 6 extensiones de DQN \\
\bottomrule
\end{tabular}
\end{table*}

\subsection{Definición Formal del Problema}

El problema de control en entornos complejos se formula como un \textbf{Proceso de Decisión de Markov} (MDP):

\begin{equation}
  \mathcal{M} = \langle \mathcal{S},\, \mathcal{A},\, \mathcal{P},\, \mathcal{R},\, \gamma \rangle
\end{equation}

La función de valor-acción óptima $Q^*$ satisface las \textbf{Ecuaciones de Bellman de Optimalidad}:

\begin{equation}
  Q^*(s,a) = \mathcal{R}(s,a) + \gamma \sum_{s' \in \mathcal{S}} \mathcal{P}(s'\mid s,a) \max_{a'} Q^*(s',a')
  = (\mathcal{T}^* Q^*)(s,a)
\end{equation}

El operador $\mathcal{T}^*$ es una contracción en norma $\ell_\infty$ con factor $\gamma < 1$, garantizando la existencia y unicidad del punto fijo \cite{Bellman1957}. Cuando $|\mathcal{S}|$ es inmanejable, se busca $\hat{q}(\cdot\,;\,\boldsymbol{\theta}) \in \mathcal{F}_{\boldsymbol{\theta}}$ tal que:

\begin{equation}
  \hat{q}(\cdot\,;\,\boldsymbol{\theta}) \approx Q^*, \qquad \boldsymbol{\theta} \in \mathbb{R}^p,\; p \ll |\mathcal{S}| \cdot |\mathcal{A}|
\end{equation}

\subsection{Contexto y Antecedentes}

\subsubsection{El Entorno \texttt{CartPole-v1} como Benchmark}

El entorno \texttt{CartPole-v1} de Gymnasium \cite{Towers2023} implementa el problema clásico de control de un péndulo invertido. El estado es el vector $\mathbf{s} = (x,\,\dot{x},\,\theta,\,\dot{\theta})^\top \in \mathbb{R}^4$, con espacio de acciones $\mathcal{A} = \{0, 1\}$ (empuje izquierdo/derecho) y recompensa $R_{t+1} = +1$ por cada paso de supervivencia. CartPole es el benchmark estándar para validar métodos aproximados de RL: su estado continuo hace inviable la tabulación, pero su dinámica es suficientemente simple para converger en cientos de episodios.

\subsubsection{Estructura Modular del Proyecto}

El proyecto implementa el patrón \textbf{Learner--Policy--Agent} con separación estricta de responsabilidades. La clase \texttt{Agent} orquesta el bucle de entrenamiento de forma agnóstica al algoritmo. Para SARSA Semi-Gradiente, el módulo \texttt{TileCodingEnv} actúa como \textbf{Gymnasium Wrapper} que transforma la observación continua $\mathbf{s} \in \mathbb{R}^4$ en el vector de características binario $\boldsymbol{\phi}(\mathbf{s}) \in \{0,1\}^d$.

\subsection{Métodos Utilizados para Abordar el Problema}

Se implementan tres familias de métodos aproximados en orden creciente de complejidad:

\begin{description}
  \item[Familia 1. Aproximación Lineal con Tile Coding (SARSA Semi-Gradiente):] $\hat{q}(\mathbf{s}, a\,;\, \mathbf{W}) = \mathbf{w}_a^\top \boldsymbol{\phi}(\mathbf{s})$, con $d = 4 \cdot 8^4 = 16\,384$. Convergencia demostrada bajo condiciones de regularidad. Limitación: no puede representar interacciones no lineales.
  \item[Familia 2. Deep Q-Network (DQN):] $Q_{\boldsymbol{\theta}}: \mathbb{R}^4 \to \mathbb{R}^2$ con arquitectura $4 \to 64 \to 2$ (ReLU). Capacidad no lineal con Experience Replay para estabilidad del gradiente.
  \item[Familia 3. Double DQN:] Extensión de DQN con desacoplamiento Selección--Evaluación para reducir el sesgo de sobreestimación positivo del operador $\max$.
\end{description}

\subsection{Diferenciación respecto a Enfoques Previos}

\begin{table*}[htb!]
\caption{Comparación entre métodos tabulares (Cap.~2) y métodos aproximados (Cap.~3)}\label{tab:diff}
\centering\small
\begin{tabular}{@{}lll@{}}
\toprule
\textbf{Dimensión} & \textbf{Métodos Tabulares} & \textbf{Presente Trabajo} \\
\midrule
Espacio de estados & Discreto, finito & Continuo, $\mathcal{S} \subseteq \mathbb{R}^4$ \\
Representación de $Q$ & Tabla $\mathbf{Q} \in \mathbb{R}^{|\mathcal{S}|\times|\mathcal{A}|}$ & Función paramétrica $\hat{q}(\cdot;\boldsymbol{\theta})$ \\
Generalización & Ninguna (estado a estado) & Implícita (vecindad en $\mathcal{S}$) \\
Número de parámetros & $|\mathcal{S}| \cdot |\mathcal{A}|$ (exponencial) & $p \ll |\mathcal{S}| \cdot |\mathcal{A}|$ \\
Entorno de evaluación & \texttt{FrozenLake-v1} (16 estados) & \texttt{CartPole-v1} ($\mathcal{S} \subseteq \mathbb{R}^4$) \\
\bottomrule
\end{tabular}
\end{table*}

% ============================================================
\section{Algoritmos}
% ============================================================

Esta sección detalla los tres algoritmos implementados en el proyecto, describiendo sus pasos, ecuaciones de actualización, pseudocódigos formales y la justificación de su elección.

\subsection{Construcción del Vector de Características: Tile Coding}

Dado un estado $\mathbf{s} = (s_1, s_2, s_3, s_4) \in \mathbb{R}^4$, se definen $m$ particiones (\textit{tilings}) del espacio de estado. Para el tiling $k$ y dimensión $j$:

\begin{align}
  o_j^{(k)} &= \frac{(k-1) \cdot \Delta_j}{m}, \quad \Delta_j = \frac{s_j^{\max} - s_j^{\min}}{n_j} \\[4pt]
  b_j^{(k)}(s_j) &= \left\lfloor \frac{s_j - s_j^{\min} + o_j^{(k)}}{\Delta_j} \right\rfloor \bmod n_j \\[4pt]
  \text{idx}^{(k)}(\mathbf{s}) &= \sum_{j=1}^{4} b_j^{(k)}(s_j) \cdot \prod_{j'=1}^{j-1} n_{j'} \\[4pt]
  \boldsymbol{\phi}(\mathbf{s}) &= \bigoplus_{k=1}^{m} \mathbf{e}_{\text{idx}^{(k)}(\mathbf{s})}^{(N_{\text{tile}})} \;\in\; \{0,1\}^{m \cdot N_{\text{tile}}}
\end{align}

El vector resultante tiene exactamente $m$ componentes activas: $\|\boldsymbol{\phi}(\mathbf{s})\|_1 = m$. Los parámetros de Tile Coding se detallan en la Tabla~\ref{tab:tilecoding}.

\begin{table}[htb!]
\caption{Parámetros de Tile Coding en el experimento}\label{tab:tilecoding}
\centering\small
\setlength{\tabcolsep}{4pt}
\begin{tabular}{@{}lll@{}}
\toprule
\textbf{Parámetro} & \textbf{Símbolo} & \textbf{Valor} \\
\midrule
Número de tilings & $m$ & 4 \\
Bins por dimensión & $n_j$ & 8, $\forall j$ \\
Teselas por tiling & $N_{\text{tile}}$ & $8^4 = 4\,096$ \\
Dimensión del vector & $d = m \cdot N_{\text{tile}}$ & 16\,384 \\
Tasa normalizada & $\alpha_{\text{eff}} = \alpha / m$ & 0.025 \\
\bottomrule
\end{tabular}
\end{table}

\subsection{SARSA Semi-Gradiente con Tile Coding}

SARSA Semi-Gradiente es un método \textbf{on-policy TD(0)} con aproximación lineal. La regla de actualización es:

\begin{equation}
  \mathbf{w}_{a,t+1} \leftarrow \mathbf{w}_{a,t} + \alpha_{\text{eff}}\,\delta_t\,\boldsymbol{\phi}(\mathbf{s}_t)
\end{equation}

con el error TD on-policy:

\begin{equation}
  \delta_t = R_{t+1} + \gamma\,\mathbf{w}_{A_{t+1}}^\top\boldsymbol{\phi}(\mathbf{s}_{t+1}) - \mathbf{w}_{A_t}^\top\boldsymbol{\phi}(\mathbf{s}_t)
\end{equation}

En forma matricial compacta (actualización de rango 1 sobre $\mathbf{W} \in \mathbb{R}^{|\mathcal{A}| \times d}$):

\begin{equation}
  \mathbf{W}_{t+1} = \mathbf{W}_t + \alpha_{\text{eff}}\,\delta_t\,\mathbf{e}_{A_t}\,\boldsymbol{\phi}(\mathbf{s}_t)^\top
\end{equation}

\begin{algorithm}
\caption{SARSA Semi-Gradiente (on-policy, TD(0))}\label{algo:sarsa}
\begin{algorithmic}[1]
\Require env, $m{=}4$, $\alpha{=}0.1$, $\gamma{=}0.95$, $\varepsilon{=}0.1$, $T_{\text{ep}}{=}400$
\Ensure $\mathbf{W} \in \mathbb{R}^{|\mathcal{A}|\times d}$ (pesos aprendidos)
\State $\mathbf{W} \gets \mathbf{0}$;\; $\alpha_{\text{eff}} \gets \alpha/m$
\For{episodio $e = 1, \ldots, T_{\text{ep}}$}
  \State $s \gets \texttt{env.reset()}$ \Comment{$s = \boldsymbol{\phi}(s_{\text{raw}})$ vía TileCodingEnv}
  \State $A \gets \varepsilon\text{-greedy}(\mathbf{W} \cdot s)$
  \While{\textbf{not} done}
    \State $s', R, \text{done} \gets \texttt{env.step}(A)$
    \State $A' \gets \varepsilon\text{-greedy}(\mathbf{W} \cdot s')$
    \State $\delta \gets R + \gamma\,(\mathbf{W}[A'] \cdot s') - (\mathbf{W}[A] \cdot s)$
    \State $\mathbf{W}[A] \gets \mathbf{W}[A] + \alpha_{\text{eff}}\cdot\delta\cdot s$
    \State $s \gets s'$;\; $A \gets A'$
  \EndWhile
\EndFor
\State \Return $\mathbf{W}$
\end{algorithmic}
\end{algorithm}

\subsection{Deep Q-Network (DQN)}

DQN \cite{Mnih2015} extiende Q-Learning a aproximadores no lineales mediante Experience Replay y red target. La arquitectura es $Q_{\boldsymbol{\theta}}: \mathbb{R}^4 \to \mathbb{R}^2$ con $|\boldsymbol{\theta}| = 450$ parámetros:

\begin{align}
  \mathbf{h}^{(1)} &= \text{ReLU}\!\left(\mathbf{W}^{(1)}\mathbf{s} + \mathbf{b}^{(1)}\right), \quad \mathbf{W}^{(1)} \in \mathbb{R}^{64\times4} \\
  Q_{\boldsymbol{\theta}}(\mathbf{s}) &= \mathbf{W}^{(2)}\mathbf{h}^{(1)} + \mathbf{b}^{(2)}, \quad \mathbf{W}^{(2)} \in \mathbb{R}^{2\times64}
\end{align}

La función de pérdida MSE sobre el mini-batch $\mathcal{B} \overset{\text{i.i.d.}}{\sim} \mathcal{D}$:

\begin{equation}
  \mathcal{L}(\boldsymbol{\theta};\,\mathcal{B}) = \frac{1}{B}\sum_{i=1}^{B}\left(y_i - Q_{\boldsymbol{\theta}}(s_i)_{a_i}\right)^2, \qquad y_i = r_i + \gamma\,(1-d_i)\,\max_{a'} Q_{\boldsymbol{\theta}}(s'_i)_{a'}
\end{equation}

Los pesos se actualizan con el optimizador \textbf{Adam} \cite{Kingma2015} ($\eta = 10^{-3}$, $\beta_1 = 0.9$, $\beta_2 = 0.999$).

\begin{algorithm}
\caption{Deep Q-Network (DQN)}\label{algo:dqn}
\begin{algorithmic}[1]
\Require env, $\eta{=}10^{-3}$, $\gamma{=}0.95$, $\varepsilon{=}0.1$, $B{=}64$, $|\mathcal{D}|_{\max}{=}10000$, $T_{\text{ep}}{=}400$
\Ensure $\boldsymbol{\theta}$ (pesos de la Q-red)
\State $\boldsymbol{\theta} \gets \text{Xavier-init}(4{\to}64{\to}2)$;\; $\mathcal{D} \gets \emptyset$
\For{episodio $e = 1, \ldots, T_{\text{ep}}$}
  \State $s \gets \texttt{env.reset()}$
  \While{\textbf{not} done}
    \State $a \gets \varepsilon\text{-greedy}(Q_{\boldsymbol{\theta}}(s))$
    \State $s', r, \text{done} \gets \texttt{env.step}(a)$
    \State $\mathcal{D} \gets \mathcal{D} \cup \{(s, a, r, s', \text{done})\}$
    \If{$|\mathcal{D}| \geq B$}
      \State $\mathcal{B} \gets \text{SampleUniform}(\mathcal{D}, B)$
      \State $y_i \gets r_i + \gamma(1{-}d_i)\max_{a'} Q_{\boldsymbol{\theta}}(s'_i)_{a'}$ \textbf{for each} $(s_i,a_i,r_i,s'_i,d_i) \in \mathcal{B}$
      \State $\boldsymbol{\theta} \gets \text{Adam}(\boldsymbol{\theta},\,\nabla_{\boldsymbol{\theta}}\mathcal{L})$
    \EndIf
    \State $s \gets s'$
  \EndWhile
\EndFor
\State \Return $\boldsymbol{\theta}$
\end{algorithmic}
\end{algorithm}

\subsection{Double Deep Q-Network (Double DQN)}

Double DQN \cite{vanHasselt2016} corrige el sesgo de sobreestimación desacoplando la selección y evaluación de la acción óptima mediante la \textbf{red target} $\boldsymbol{\theta}^-$:

\begin{equation}
  y_i^{\text{DDQN}} = r_i + \gamma\,(1-d_i)\; Q_{\boldsymbol{\theta}^-}\!\!\left(s'_i,\; \underbrace{\arg\max_{a'} Q_{\boldsymbol{\theta}}(s'_i, a')}_{\text{selección: red online } \boldsymbol{\theta}}\right)
\end{equation}

La red target se actualiza con copia \textit{hard} cada $\tau = 100$ pasos: $\boldsymbol{\theta}^- \leftarrow \boldsymbol{\theta}$.

\begin{algorithm}
\caption{Double Deep Q-Network (Double DQN)}\label{algo:ddqn}
\begin{algorithmic}[1]
\Require env, $\eta{=}10^{-3}$, $\gamma{=}0.95$, $\varepsilon{=}0.1$, $B{=}64$, $\tau{=}100$, $T_{\text{ep}}{=}400$
\Ensure $\boldsymbol{\theta}$ (pesos de la red online)
\State $\boldsymbol{\theta} \gets \text{Xavier-init}(4{\to}64{\to}2)$;\; $\boldsymbol{\theta}^- \gets \boldsymbol{\theta}$;\; $\mathcal{D} \gets \emptyset$;\; $t_g \gets 0$
\For{episodio $e = 1, \ldots, T_{\text{ep}}$}
  \State $s \gets \texttt{env.reset()}$
  \While{\textbf{not} done}
    \State $a \gets \varepsilon\text{-greedy}(Q_{\boldsymbol{\theta}}(s))$
    \State $s', r, \text{done} \gets \texttt{env.step}(a)$
    \State $\mathcal{D} \gets \mathcal{D} \cup \{(s, a, r, s', \text{done})\}$
    \If{$|\mathcal{D}| \geq B$}
      \State $\mathcal{B} \gets \text{SampleUniform}(\mathcal{D}, B)$
      \State $a^*_i \gets \arg\max_{a'} Q_{\boldsymbol{\theta}}(s'_i)_{a'}$;\; $y_i \gets r_i + \gamma(1{-}d_i)\,Q_{\boldsymbol{\theta}^-}(s'_i)_{a^*_i}$
      \State $\boldsymbol{\theta} \gets \text{Adam}(\boldsymbol{\theta},\,\nabla_{\boldsymbol{\theta}}\mathcal{L})$
    \EndIf
    \State $t_g \gets t_g + 1$
    \If{$t_g \bmod \tau = 0$} $\boldsymbol{\theta}^- \gets \boldsymbol{\theta}$ \EndIf
    \State $s \gets s'$
  \EndWhile
\EndFor
\State \Return $\boldsymbol{\theta}$
\end{algorithmic}
\end{algorithm}

\subsection{Justificación de los Algoritmos Elegidos}

\begin{table*}[htb!]
\caption{Criterios de selección de los algoritmos}\label{tab:criterios}
\centering\small
\begin{tabular}{@{}lcccc@{}}
\toprule
\textbf{Criterio} & \textbf{SARSA SG} & \textbf{DQN} & \textbf{DDQN} \\
\midrule
Espacio continuo $\mathcal{S} \subseteq \mathbb{R}^4$ & \checkmark & \checkmark & \checkmark \\
Convergencia demostrada & \checkmark & Parcial & Parcial \\
Corrección de sobreestimación & --- & --- & \checkmark \\
Generalización no lineal & --- & \checkmark & \checkmark \\
Eficiencia muestral & Alta ($O(m)$) & Media & Media \\
\bottomrule
\end{tabular}
\end{table*}

SARSA Semi-Gradiente actúa como \textbf{línea base} clásica. DQN es el representante canónico del RL profundo basado en valor. Double DQN aísla empíricamente el efecto de la corrección del sesgo de sobreestimación, al compartir arquitectura y todos los hiperparámetros con DQN. Los criterios de selección se resumen en la Tabla~\ref{tab:criterios}.

% ============================================================
\section{Evaluación / Experimentos}
% ============================================================

\subsection{Configuración Experimental: herramientas y entorno}

Los experimentos se desarrollaron sobre el siguiente stack tecnológico, recogido en la Tabla~\ref{tab:stack}:

\begin{table*}[htb!]
\caption{Stack tecnológico del experimento}\label{tab:stack}
\centering\small
\begin{tabular}{@{}llp{7cm}@{}}
\toprule
\textbf{Componente} & \textbf{Herramienta} & \textbf{Función} \\
\midrule
Entorno RL & \texttt{gymnasium} \cite{Towers2023} & Simulador CartPole-v1 \\
Álgebra vectorial & \texttt{numpy} & Tile Coding, actualización de pesos \\
Deep Learning & \texttt{torch} (PyTorch) & Redes neuronales, Adam \\
Visualización & \texttt{matplotlib} & Curvas de aprendizaje \\
Monitoreo & \texttt{tqdm} & Progreso del entrenamiento \\
Buffer replay & \texttt{collections.deque} & Experience Replay \\
\bottomrule
\end{tabular}
\end{table*}

\subsubsection{Entorno de Prueba: \texttt{CartPole-v1}}

El entorno \texttt{CartPole-v1} induce un MDP con espacio de estados continuo de dimensión 4:

\begin{equation}
  \mathbf{s} = \begin{pmatrix} x \\ \dot{x} \\ \theta \\ \dot{\theta} \end{pmatrix} \in \mathbb{R}^4
\end{equation}

\begin{table}[htb!]
\caption{Variables de estado de \texttt{CartPole-v1}}\label{tab:cartpole}
\centering\small
\setlength{\tabcolsep}{4pt}
\begin{tabular}{@{}lll@{}}
\toprule
\textbf{Variable} & \textbf{Descripción} & \textbf{Rango} \\
\midrule
$x$ & Posición del carro & $[-4.8,\ 4.8]$ \\
$\dot{x}$ & Velocidad lineal & $\mathbb{R}$ \\
$\theta$ & Ángulo del palo (rad) & $(-\pi/15,\ \pi/15)$ \\
$\dot{\theta}$ & Velocidad angular & $\mathbb{R}$ \\
\bottomrule
\end{tabular}
\end{table}

Las acciones son $\mathcal{A} = \{0\ (\text{izq.}),\ 1\ (\text{der.})\}$ y la recompensa es $R_{t+1} = +1$ por cada paso de supervivencia (máximo 500 por episodio) Las variables de estado de \texttt{CartPole-v1} se detallan en la Tabla~\ref{tab:cartpole}.

\subsubsection{Protocolo Experimental}

\begin{table}[htb!]
\caption{Protocolo experimental}\label{tab:protocolo}
\centering
\begin{tabular}{@{}ll@{}}
\toprule
\textbf{Parámetro} & \textbf{Valor} \\
\midrule
Número de ejecuciones ($N_\text{runs}$) & 10 \\
Episodios por ejecución ($T_\text{ep}$) & 400 \\
Semilla base ($\sigma$) & 123 \\
Factor de descuento ($\gamma$) & 0.95 \\
Política de exploración & $\varepsilon$-greedy, $\varepsilon = 0.1$ \\
\bottomrule
\end{tabular}
\end{table}

El protocolo experimental se detalla en la tabla \ref{tab:protocolo}. La reproducibilidad se garantiza mediante la función \texttt{set\_seed}($\sigma$).


\subsubsection{Datasets Empleados}

Por tratarse de un entorno de simulación RL, no existe un dataset estático previo. Las muestras son \textbf{trayectorias} $(s_t, a_t, r_{t+1}, s_{t+1})$ generadas en línea. Los métodos Deep RL almacenan transiciones en el Experience Replay Buffer $\mathcal{D}$ ($|\mathcal{D}|_\max = 10\,000$) y extraen mini-batches $\mathcal{B}$ de tamaño $B = 64$:

\begin{equation}
  \mathcal{B} = \{(s_i, a_i, r_i, s'_i, d_i)\}_{i=1}^{64} \overset{\text{i.i.d.}}{\sim} \mathcal{D}
\end{equation}

\subsection{Métricas de Evaluación}

\subsubsection{Métrica Principal: $\overline{R}_{[-50:]}$}

La métrica principal es la recompensa media en los últimos 50 episodios del entrenamiento, promediada sobre los 10 runs:

\begin{equation}
  \overline{R}_{[-50:]} = \frac{1}{50} \sum_{e=351}^{400} \bar{R}_e, \qquad \bar{R}_e = \frac{1}{N_\text{runs}} \sum_{n=1}^{N_\text{runs}} R_e^{(n)}
\end{equation}

Un valor de 500 indica política perfecta (palo siempre erguido).

\subsubsection{Pérdida de Entrenamiento (Solo Deep RL)}

Para DQN y Double DQN se registra la pérdida media por episodio:

\begin{equation}
  \bar{\mathcal{L}}_e = \frac{1}{|\text{updates}_e|} \sum_{t \in \text{updates}_e} \mathcal{L}(\boldsymbol{\theta}_t;\, \mathcal{B}_t)
\end{equation}

\begin{table*}[htb!]
\caption{Resumen de métricas de evaluación}\label{tab:metricas}
\centering\small
\begin{tabular}{@{}llll@{}}
\toprule
\textbf{Métrica} & \textbf{Símbolo} & \textbf{Propósito} & \textbf{Algoritmos} \\
\midrule
Recompensa régimen estacionario & $\overline{R}_{[-50:]}$ & Calidad de la política & Todos \\
Curva de aprendizaje & $\bar{R}_e$ & Dinámica de convergencia & Todos \\
Pérdida de entrenamiento & $\bar{\mathcal{L}}_e$ & Estabilidad de optimización & DQN, DDQN \\
\bottomrule
\end{tabular}
\end{table*}

Véase Tabla~\ref{tab:metricas}.

\subsection{Resultados Obtenidos}

\begin{table*}[htb!]
\caption{Resultados globales — Recompensa media en los últimos 50 episodios}\label{tab:resultados}
\centering\small
\begin{tabular}{@{}llccc@{}}
\toprule
\textbf{Algoritmo} & \textbf{Aproximador} & \textbf{Parámetros} & $\overline{R}_{[-50:]}$ & \textbf{Relat. al máx. (500)} \\
\midrule
SARSA Semi-Gradiente & Lineal + Tile Coding & 32\,768 & 69.42 & 13.9\% \\
DQN & Red neuronal (PyTorch) & 450 & \textbf{218.26} & \textbf{43.7\%} \\
Double DQN & Dual red neuronal & 900 & 189.06 & 37.8\% \\
\bottomrule
\end{tabular}
\end{table*}

Ningún algoritmo alcanza la recompensa máxima de 500 en el horizonte de 400 episodios, indicando convergencia incompleta. El ordenamiento DQN $>$ Double DQN $>$ SARSA es consistente con las predicciones teóricas. Véase Tabla~\ref{tab:resultados}.

\subsection{Análisis Crítico y Comparación con Otros Enfoques}

\subsubsection{SARSA Semi-Gradiente — $\overline{R}_{[-50:]} = 69.42$}

La restricción de la hipótesis lineal impide capturar las interacciones no lineales entre las variables de estado (posición, velocidad, ángulo, velocidad angular correlacionan de forma compleja en CartPole). Adicionalmente, los 32\,768 pesos reciben actualizaciones que solo afectan a $m = 4$ entradas en cada paso, mientras que DQN actualiza densamente sus 450 parámetros vía backpropagation.

\subsubsection{DQN — $\overline{R}_{[-50:]} = 218.26$ (mejor resultado)}

La no-linealidad de ReLU permite representar $Q^*$ con precisión suficiente en 400 episodios. El Experience Replay estabiliza el gradiente estocástico al decorrelacionar las muestras. La principal limitación teórica es el sesgo de sobreestimación positivo acumulado en la propagación de Bellman.

\subsubsection{Double DQN — $\overline{R}_{[-50:]} = 189.06$}

El desacoplamiento Selección--Evaluación mediante la red target reduce el sesgo de sobreestimación. Sin embargo, la \textbf{inercia de la red target} ($\tau = 100$ pasos) ralentiza la convergencia inicial en el horizonte de 400 episodios. Con horizontes más largos ($T_\text{ep} > 1\,000$), la reducción del sesgo favorece estadísticamente a Double DQN, como documenta la literatura \cite{vanHasselt2016}.

En la Tabla~\ref{tab:literatura}, tenemos una comparación con resultados de la literatura.

\begin{table*}[htb!]
\caption{Comparación con resultados de la literatura}\label{tab:literatura}
\centering\small
\begin{tabular}{@{}lllll@{}}
\toprule
\textbf{Referencia} & \textbf{Algoritmo} & \textbf{Entorno} & $\overline{R}$ & \textbf{Observación} \\
\midrule
Mnih et al.~\cite{Mnih2015} & DQN & Atari & Nivel humano & $\sim$10M steps \\
van Hasselt et al.~\cite{vanHasselt2016} & DDQN & Atari & $>$ DQN en 43/49 & Horizonte largo \\
Sutton \& Barto \cite{Sutton2018} & SARSA SG & MountainCar & $-100$ a $-150$ & Comparablemente \\
\textbf{Este trabajo} & Todos & CartPole-v1 & 69/218/189 & 400~ep., incompleto \\
\bottomrule
\end{tabular}
\end{table*}

% ============================================================
\section{Conclusiones}
% ============================================================

Esta sección sintetiza los hallazgos principales del estudio experimental sobre métodos aproximados de Aprendizaje por Refuerzo, extrae las lecciones teóricas y prácticas derivadas de los resultados, discute las limitaciones del trabajo y propone líneas de investigación futuras.

\subsection{Resumen de los Principales Resultados Obtenidos}

El experimento comparó tres familias de métodos aproximados sobre el entorno \texttt{CartPole-v1} con un protocolo de $N_\text{runs} = 10$ ejecuciones y $T_\text{ep} = 400$ episodios. Los resultados evidencian un ordenamiento consistente con las predicciones teóricas: DQN $>$ Double DQN $>$ SARSA Semi-Gradiente (véase Tabla~\ref{tab:resultados}).

Se extraen cinco conclusiones clave:

\begin{enumerate}
  \item \textbf{Los métodos aproximados generalizan en espacios continuos.} La sustitución de la tabla $\mathbf{Q}$ por un aproximador paramétrico $\hat{q}(\cdot;\boldsymbol{\theta})$ permite aprender en $\mathcal{S} \subseteq \mathbb{R}^4$, imposible con representación tabular.

  \item \textbf{La capacidad del aproximador determina la calidad de la política.} La brecha SARSA (69.42) vs.\ DQN (218.26) refleja la diferencia en familia de hipótesis: la combinación lineal de tiles no puede representar las interacciones no lineales que sí captura ReLU.

  \item \textbf{El Experience Replay es condición crítica de estabilidad.} El muestreo uniforme $\mathcal{B} \overset{\text{i.i.d.}}{\sim} \mathcal{D}$ convierte el problema en regresión supervisada, eliminando el sesgo de la correlación temporal entre transiciones consecutivas.

  \item \textbf{La red target regulariza el problema no estacionario de los targets.} La copia periódica $\boldsymbol{\theta}^- \leftarrow \boldsymbol{\theta}$ cada $\tau = 100$ pasos congela la distribución de targets, convirtiendo la optimización no estacionaria en una quasi-estacionaria localmente.

  \item \textbf{El ordenamiento empírico depende del horizonte.} La ventaja teórica de Double DQN sobre DQN (sesgo de sobreestimación reducido) solo se manifiesta plenamente para $T_\text{ep} > 1\,000$; en 400 episodios, la inercia de la red target ralentiza la convergencia inicial.
\end{enumerate}

\subsection{Limitaciones del Estudio}

A pesar de los resultados satisfactorios, el presente trabajo presenta las siguientes limitaciones que deben tenerse en cuenta al interpretar las conclusiones tal como muestra la tabla \ref{tab:limitaciones}

\begin{table*}[htb!]
\caption{Limitaciones del estudio experimental}\label{tab:limitaciones}
\centering\small
\begin{tabular}{@{}lp{9cm}@{}}
\toprule
\textbf{Limitación} & \textbf{Impacto} \\
\midrule
Horizonte reducido ($T_\text{ep} = 400$) & Convergencia incompleta; ningún algoritmo alcanza 500 \\
Entorno único (\texttt{CartPole-v1}) & Rankings relativos no generalizables a otros entornos \\
Hiperparámetros fijos (sin búsqueda) & Resultados comparativos, no óptimos absolutos \\
Sin intervalos de confianza formales & No se cuantifica la varianza entre runs estadísticamente \\
Política $\varepsilon$ constante ($\varepsilon = 0.1$) & No se explota la política aprendida en fases tardías \\
\bottomrule
\end{tabular}
\end{table*}

\subsection{Líneas de Trabajo Futuro}

\begin{enumerate}
  \item \textbf{Mayor horizonte de entrenamiento} ($T_\text{ep} > 1\,000$) para confirmar la superioridad de Double DQN sobre DQN.
  \item \textbf{Decaimiento de $\varepsilon$} (\textit{epsilon decay}): $\varepsilon(t) = \varepsilon_{\min} + (\varepsilon_0 - \varepsilon_{\min}) e^{-\lambda t}$ para reducir exploración progresivamente.
  \item \textbf{Optimización bayesiana de hiperparámetros} mediante \textit{Optuna} o \textit{Ax}.
  \item \textbf{Entornos de mayor complejidad}: LunarLander-v2 ($\mathbb{R}^8$), BipedalWalker ($\mathbb{R}^{24}$), MuJoCo.
  \item \textbf{Arquitecturas avanzadas}: Dueling DQN \cite{Wang2016}, Prioritized Experience Replay \cite{Schaul2016}, Rainbow \cite{Hessel2018}.
  \item \textbf{Análisis estadístico formal}: \textit{bootstrapped confidence intervals} al 95\%, test de Wilcoxon con corrección de Bonferroni.
  \item \textbf{Transferencia de aprendizaje}: uso de los pesos $\boldsymbol{\theta}$ pre-entrenados como inicialización en entornos relacionados.
\end{enumerate}

\subsection{Reflexión Final}

Este capítulo demuestra que la transición de los métodos tabulares a los métodos aproximados no es meramente una extensión técnica, sino un cambio cualitativo en la naturaleza del problema de aprendizaje: de la memorización de una tabla finita a la \textbf{inferencia de una función} sobre un espacio continuo. Los tres algoritmos estudiados representan tres puntos del espectro de complejidad del aproximador ---lineal, red \textit{shallow}, red con target estabilizada--- y sus diferencias de rendimiento ilustran de forma concreta el compromiso entre \textbf{capacidad expresiva}, \textbf{eficiencia muestral} y \textbf{estabilidad del entrenamiento} central en la investigación moderna de Aprendizaje por Refuerzo Profundo.


% Bibliografía  
\begin{thebibliography}{9} % El parámetro 9 indica el ancho máximo de la etiqueta, aquí "9" es un ejemplo
	\setlength{\itemsep}{0pt} % No tocar. Reduce aún más el espacio entre ítems 
  	\setlength{\parsep}{0pt}    % No tocar. Elimina el espacio adicional entre párrafos dentro de un ítem
  	
    \bibitem{sutton2018reinforcement} % Para citar usar \cite{ref1}
    Sutton, R.S. y Barto, A.G (2018). \textit{Reinforcement Learning: An Introduction}. The MIT Press Cambridge, Second Edition.

    \bibitem{auer2002finitetime} % Para citar usar \cite{ref2}
    Auer, P., Cesa-Bianchi, N. y Fischer, P (2002). \textit{Finite-time Analysis of the Multiarmed Bandit Problem}. Machine Learning 47, 235–256.

    \bibitem{Watkins1989}
    Watkins, Christopher y Dayan, Peter. (1992). Technical Note: Q-Learning. Machine Learning. 8. 279-292. 10.1007/BF00992698. 

    \bibitem{SARSA1994}
    Rummery, G. & Niranjan, Mahesan. (1994). On-Line Q-Learning Using Connectionist Systems. Technical Report CUED/F-INFENG/TR 166. 

    \bibitem{Hasselt2010DD}
    Van Hasselt, Hado. (2010). Double Q-learning. 2613-2621.

    \bibitem{Mnih2015}
    Mnih, V. et al. (2015). Human-level control through deep reinforcement learning. \textit{Nature}, 518, 529--533.

    \bibitem{vanHasselt2016}
    van Hasselt, H., Guez, A. \& Silver, D. (2016). Deep Reinforcement Learning with Double Q-learning. \textit{AAAI-16}.

    \bibitem{Sutton2018}
    Sutton, R.S. \& Barto, A.G. (2018). \textit{Reinforcement Learning: An Introduction} (2nd ed.). MIT Press.

    \bibitem{Sutton1996}
    Sutton, R.S. (1996). Generalization in reinforcement learning. \textit{NIPS 1995}.

    \bibitem{Towers2023}
    Towers, M. et al. (2023). Gymnasium. \textit{Zenodo}. \url{https://doi.org/10.5281/zenodo.8127025}

    \bibitem{Kingma2015}
    Kingma, D.P. \& Ba, J. (2015). Adam: A Method for Stochastic Optimization. \textit{ICLR 2015}.

    \bibitem{Lin1992}
    Lin, L.-J. (1992). Self-improving reactive agents based on reinforcement learning. \textit{Machine Learning}, 8(3), 293--321.

    \bibitem{Bellman1957}
    Bellman, R. (1957). \textit{Dynamic Programming}. Princeton University Press.

    \bibitem{Silver2017}
    Silver, D. et al. (2017). Mastering the game of Go without human knowledge. \textit{Nature}, 550, 354--359.

    \bibitem{Wang2016}
    Wang, Z. et al. (2016). Dueling Network Architectures for Deep Reinforcement Learning. \textit{ICML 2016}.

    \bibitem{Schaul2016}
    Schaul, T. et al. (2016). Prioritized Experience Replay. \textit{ICLR 2016}.

    \bibitem{Hessel2018}
    Hessel, M. et al. (2018). Rainbow: Combining Improvements in Deep Reinforcement Learning. \textit{AAAI-18}.
    
\end{thebibliography}


Listar todas las referencias utilizadas en el trabajo, siguiendo un formato bibliográfico estándar de \LaTeX\/ o de alguna asociación (APA, IEEE, etc.). La información obligatoria para que tipo de publicación la puede consultar en \url{https://bibtex.eu/fields/} Se recomienda utilizar un gestor de referencias como Bib\TeX\/ para facilitar la organización de las citas solo si va a poner mucha bibliografía. 


\end{document} 